\graphicspath{{Imagenes/}}

% ============================================================
%  SECCIÓN 1 – Introducción y Motivación
% ============================================================

\section{Introducción y Motivación}

\subsection{Contexto industrial: empresa AUTOELEC}

La empresa \textbf{AUTOELEC} se dedica al asesoramiento, venta, reparación y alquiler de autoelevadores y baterías. Dentro de su proceso productivo, una tarea crítica es el \textbf{ensamble de paquetes de baterías de plomo-ácido} destinadas a autoelevadores industriales. Estos paquetes se construyen ubicando manualmente múltiples vasos (celdas individuales) dentro de un rack metálico, conformando la batería final que alimenta el vehículo.

\begin{figure}[H]
    \centering
    \begin{minipage}[t]{0.47\textwidth}
        \centering
        \includegraphics[width=\linewidth]{Celdas_Carga.png}
        \caption{Vasos de plomo-ácido dispuestos en rack previo al ensamble.}
        \label{fig:celdas_carga}
    \end{minipage}
    \hfill
    \begin{minipage}[t]{0.47\textwidth}
        \centering
        \includegraphics[width=\linewidth]{Bateria.png}
        \caption{Paquete de baterías terminado, marca Eternity -- distribuido por AUTOELEC.}
        \label{fig:bateria}
    \end{minipage}
\end{figure}

\subsection{Descripción del problema}

El proceso de ensamble de vasos, cuyo peso varía entre 20 y 100\,kg, requiere un control de movimientos eficiente y una capacidad operativa que minimice errores, reduzca los tiempos de producción y evite daños derivados de una manipulación inadecuada de los materiales. Actualmente, al tratarse de un proceso \textbf{completamente manual}, se generan riesgos ergonómicos que provocan un desgaste prematuro en la salud del operador, quien está expuesto a esfuerzos físicos considerables y repetitivos.

El operador trabaja con un rack destinado al montaje de múltiples vasos de plomo-ácido. Para efectuar el ensamble, debe elevar manualmente cada vaso ---asistido únicamente por su propia fuerza física--- y ubicarlo en una posición determinada dentro del rack. Una vez instaladas la mayoría de las unidades, el espacio libre en el rack resulta muy reducido, lo que dificulta las maniobras de inserción y obliga al operador a trabajar en posiciones forzadas que pueden causar accidentes.

\subsection{Solución propuesta en Robótica II: manipulador electrohidráulico}

Para abordar esta problemática, en la asignatura Robótica II se propuso el desarrollo de un \textbf{Manipulador de Asistencia Ergonómica} de tipo electrohidráulico. Este concepto soporta el peso de la carga y permite al operador guiar el movimiento con un esfuerzo mínimo, siguiendo el principio de los llamados \textit{manipuladores industriales ingrávidos}.

De las alternativas evaluadas (neumático, eléctrico y contrapeso), la opción seleccionada fue un \textbf{brazo articulado con contrapeso y accionamiento electrohidráulico} de simple efecto. En esta configuración:

\begin{itemize}
    \item El contrapeso equilibra estáticamente el extremo distal del brazo.
    \item Un motor eléctrico acciona una bomba hidráulica que eleva la carga a velocidad constante al presionar el pulsador de elevación en la empuñadura.
    \item Al activar el pulsador de descenso, el retorno controlado del fluido hidráulico baja la carga por gravedad a velocidad constante.
    \item El control integrado en la empuñadura ergonómica permite que el operador trabaje adyacente a la pieza durante todo el ciclo de ensamble, facilitando el enganche, desenganche y posicionamiento final.
\end{itemize}

Este diseño permite \textbf{acelerar significativamente los ciclos de manipulación} en comparación con el proceso manual, reduciendo a la vez la fatiga física y el riesgo de lesiones del operador.

\subsection{Motivación para el entorno de Realidad Virtual}

La teleoperación y el aprendizaje de sistemas de manipulación robótica presentan desafíos particulares: existen \textit{riesgos operativos} durante la fase de entrenamiento (colisiones con el rack, caída de vasos, golpes al operador), se requiere \textit{precisión manual} en el posicionamiento de piezas pesadas y hay \textit{alta variabilidad} en la destreza inicial de distintos operadores.

Un entorno de \textbf{Realidad Virtual (VR)} permite abordar estos desafíos de forma integral:

\begin{itemize}
    \item \textbf{Seguridad}: el operador aprende a manejar el manipulador sin riesgo de dañar equipos ni lastimarse, al interactuar con una réplica virtual del sistema real.
    \item \textbf{Reducción de costos}: se eliminan los costos asociados a errores durante el entrenamiento (materiales dañados, paradas de producción).
    \item \textbf{Iteración rápida}: el entorno virtual puede reconfigurarse para distintas secuencias de ensamble, cargas o restricciones sin intervención física.
    \item \textbf{Métricas objetivas}: es posible instrumentar automáticamente indicadores de desempeño como tiempo por tarea, suavidad del movimiento (\textit{jerk}), colisiones y reintentos.
    \item \textbf{Transferencia de expertise}: los patrones de movimiento aprendidos en VR pueden transferirse al control del robot real a través del mismo joystick físico utilizado durante el entrenamiento.
\end{itemize}

En este contexto, el presente proyecto propone el desarrollo de un \textbf{entorno VR para entrenar operarios} en la teleoperación del manipulador electrohidráulico de AUTOELEC, utilizando un joystick diseñado ad hoc sobre \textit{Raspberry Pi Pico} como interfaz de control, y el motor gráfico \textit{Unity} como plataforma de simulación y experiencia inmersiva.

% ============================================================
%  SECCIÓN 2 – Objetivos y Alcance
% ============================================================

\section{Objetivos y Alcance}

\subsection{Objetivo general}

Desarrollar un entorno VR y un esquema de control que permitan entrenar y luego teleoperar un brazo robótico mediante un joystick propio, priorizando seguridad, usabilidad y desempeño.

\subsection{Objetivos específicos}

\begin{enumerate}
    \item Implementar un \textbf{simulador} de brazo robótico electrohidráulico en Unity con interacción XR.
    \item Diseñar el \textbf{mapeo de controles} del joystick (RPi Pico) a los grados de libertad del robot.
    \item Instrumentar \textbf{métricas de entrenamiento}: tiempo por tarea, errores, colisiones, suavidad/jerk.
    \item Diseñar y validar \textbf{modos de teleoperación}: velocidad limitada, asistencia, seguimiento de trayectoria.
    \item Incorporar \textbf{mecanismos de seguridad}: parada de emergencia (\textit{e-stop}), zonas prohibidas, límites de torque/velocidad.
    \item Realizar una \textbf{evaluación} con usuarios: usabilidad (SUS), curva de aprendizaje y transferencia al robot real.
\end{enumerate}

\subsection{Alcance}

\begin{itemize}
    \item Entrenamiento VR de tareas representativas de la aplicación (ej.: \textit{pick \& place} de vasos, posicionamiento preciso dentro del rack).
    \item Teleoperación con joystick RPi Pico (prototipo funcional de hardware y firmware).
    \item Integración de comunicación (serial/UDP/TCP) entre el joystick y Unity.
    \item Capa de seguridad básica en la simulación.
\end{itemize}

\subsection{Fuera de alcance (versión actual)}

\begin{itemize}
    \item Retroalimentación háptica avanzada (vibración proporcional a la fuerza).
    \item \textit{Motion planning} autónomo o visión por computadora.
    \item Integración física directa con el robot electrohidráulico real (queda como trabajo futuro).
    \item Certificación normativa industrial.
\end{itemize}

% ============================================================
%  SECCIÓN 3 – Metodología y Tecnologías
% ============================================================

\section{Metodología y Tecnologías}

\subsection{Motor gráfico y entorno XR}

El proyecto utiliza \textbf{Unity 6.3 LTS} como motor gráfico principal. Unity provee soporte nativo para \textit{Extended Reality} (XR) a través del paquete \textit{XR Interaction Toolkit}, así como un sistema de \textit{input} flexible (Unity Input System) que permite integrar dispositivos HID personalizados como el joystick RPi Pico. Se emplea el pipeline de renderizado Universal Render Pipeline (URP), compatible con los principales \textit{headsets} de VR del mercado y optimizado para rendimiento en tiempo real.

\subsection{Arquitectura del sistema}

El sistema se organiza en capas funcionales que delimitan las responsabilidades de cada componente:

\begin{itemize}
    \item \textbf{Capa de Dispositivos}: el joystick RPi Pico se conecta al PC mediante USB, exponiendo sus ejes analógicos y botones como dispositivo HID estándar.
    \item \textbf{Capa de Control}: módulo en Unity responsable del mapeo de los ejes del joystick a comandos de velocidad o posición para los grados de libertad del robot.
    \item \textbf{Capa de Simulación}: modelo cinemático del brazo robótico, físicas de colisión y restricciones articulares dentro del entorno Unity.
    \item \textbf{Capa de VR/UX}: escenas de entrenamiento, HUD de retroalimentación, objetivos interactivos y tutoriales progresivos.
    \item \textbf{Capa de Teleoperación}: canal de comunicación serie/UDP y \textit{safety layer} para la eventual conexión con el robot real.
\end{itemize}

\subsection{Hardware de interfaz: Joystick RPi Pico}

El hardware de control fue diseñado e implementado íntegramente por Francisco Castel como componente específico del proyecto.

\subsubsection*{Diseño y prototipado}

Se diseñó un controlador de dos ejes analógicos con carcasa ergonómica impresa en 3D. El joystick utiliza un microcontrolador \textbf{Raspberry Pi Pico} como unidad central, aprovechando sus conversores analógico-digitales (ADC) para la lectura de los potenciómetros de cada eje y sus capacidades USB para la comunicación con el host.

\subsubsection*{Firmware y protocolo de comunicación}

El firmware fue desarrollado en C/C++ sobre el SDK de RPi Pico. Para la capa de comunicación se implementó el protocolo \textbf{USB HID} utilizando la librería \textit{TinyUSB}, lo que permite que el joystick sea reconocido automáticamente por el sistema operativo como un gamepad estándar, sin necesidad de drivers adicionales. Esto simplifica significativamente la integración con Unity a través de su Input System.

El funcionamiento correcto del protocolo HID fue verificado con la herramienta \textit{Gamepad Tester}, que confirmó la lectura de los ejes analógicos y la identificación del dispositivo como \texttt{cafe-4004-Joystick VR}.

\begin{figure}[H]
    \centering
    \includegraphics[angle=90, width=0.65\textwidth]{joystick_top_inside.png}
    \caption{Interior del joystick: RPi Pico montado sobre placa de prototipado con los módulos analógicos de cada eje.}
    \label{fig:joystick_inside}
\end{figure}

\begin{figure}[H]
    \centering
    \includegraphics[angle=90, width=0.65\textwidth]{joystick_top_full.png}
    \caption{Vista exterior del prototipo con carcasa impresa en 3D y ejes de control montados.}
    \label{fig:joystick_full}
\end{figure}

\begin{figure}[H]
    \centering
    \includegraphics[width=0.65\textwidth]{gamepad_tester.png}
    \caption{Verificación del dispositivo HID en \textit{Gamepad Tester}: lectura en tiempo real de los ejes analógicos del joystick RPi Pico.}
    \label{fig:gamepad_tester}
\end{figure}

\subsubsection*{Estado actual y próximos pasos}

El primer prototipo funcional está operativo, con firmware básico de lectura ADC y comunicación USB HID. Las iteraciones planificadas incluyen ajuste de zona muerta (\textit{deadband}), filtrado de ruido (\textit{debounce}), refinamiento ergonómico de la carcasa y evaluación de vibración háptica opcional.

\subsection{Simulación del robot en Unity}

La simulación del brazo robótico electrohidráulico se implementa dentro de Unity utilizando el paquete de robótica \textbf{Flange}, que provee modelos 3D articulados y una interfaz de control integrada en el Inspector de Unity. Este paquete expone un panel de control con sliders por cada articulación (J1--J6), posicionamiento cartesiano y funciones de guardado de pose.

La física de las articulaciones se gestiona mediante componentes \texttt{HingeJoint} y \texttt{Rigidbody} de Unity, que permiten simular las restricciones mecánicas del manipulador real. Los grados de libertad del robot son controlados mediante scripts C\# que reciben comandos de velocidad desde la capa de control del joystick.

\begin{figure}[H]
    \centering
    \includegraphics[width=0.92\textwidth]{KUKA_Flange.png}
    \caption{Escena de simulación en Unity: brazo KUKA KR210 con efector final tipo gancho integrado mediante el paquete \emph{Flange}. Al fondo se aprecia el rack de ensamble y otros elementos del entorno industrial virtual.}
    \label{fig:kuka_flange}
\end{figure}

Los detalles del desarrollo de la simulación ---incluyendo los scripts implementados, la arquitectura de control y el modo cámara--- se describen en la sección siguiente.

% ============================================================
%  SECCIÓN 4 – Sistema de Control por Joystick en Entorno Virtual
% ============================================================

\section{Sistema de Control por Joystick en Entorno Virtual}
\label{sec:control_joystick}

\subsection{Descripción general}

El sistema de control desarrollado en Unity permite operar el manipulador robótico,
accionar la garra, navegar por la escena y utilizar el menú sin necesidad del
teclado ni del ratón. \texttt{InputProfileSwitcher} permite seleccionar dos perfiles
de entrada: el control PlayStation~4 y el joystick propio VR--2 basado en RPi Pico.
Las acciones de cada perfil se asignan a \texttt{JoystickAdapter},
\texttt{CameraJoystickController} y al controlador de la garra mediante
\emph{Unity Input System}.

\subsection{Definición de las entradas del joystick}

El perfil PS4 utiliza el archivo \texttt{PS4Joystick.inputactions}. La
Tabla~\ref{tab:ps4_commands} resume el mapeo físico y todas las funciones de
operación disponibles.

\begin{table}[H]
\centering
\small
\setlength{\tabcolsep}{4pt}
\renewcommand{\arraystretch}{1.18}
\caption{Comandos implementados para el control PlayStation~4.}
\label{tab:ps4_commands}
\begin{tabular}{p{3.0cm}p{3.2cm}p{9.4cm}}
\hline
\textbf{Control físico} & \textbf{Acción de entrada} & \textbf{Función} \\
\hline
Stick izquierdo & \texttt{MoveX}, \texttt{MoveZ} & Traslada el TCP en el plano horizontal; en cámara desplaza hacia delante/atrás y lateralmente; en modo J6 interpreta el movimiento circular de la palanca. \\
Stick derecho & \texttt{MoveY}, \texttt{ViewUp}, \texttt{ViewSide} & Su eje vertical mueve el TCP en altura. En cámara, ambos ejes controlan \textit{pitch} y \textit{yaw}. \\
L3, clic corto & \texttt{ModoCamara} & Conmuta entre modo Robot y modo Cámara. El cambio se difiere \SI{0.4}{\second} para distinguirlo de un doble clic. \\
L3, doble clic & \texttt{ModoCamara} & Activa o desactiva el modo J6 exclusivo. \\
L3, mantener \SI{2}{\second} & \texttt{ModoCamara} & Abre o cierra el menú de pausa sin ejecutar el cambio de cámara. \\
R3, clic corto & \texttt{Gripper} & Alterna la garra entre abierta y cerrada al soltar el botón. \\
R3, mantener \SI{1}{\second} & \texttt{Gripper} & Ejecuta el retorno de J6 a su referencia física; no alterna la garra. \\
L1 / R1 & \texttt{J6AntiHor}, \texttt{J6Hor} & En modo J6, gira en sentido antihorario/horario a velocidad constante. \\
Cuadrado & \texttt{J6Home} & Inicia directamente el retorno de J6 a referencia. \\
Options & \texttt{Pause} & Abre o cierra el menú de pausa. \\
Cruceta / Cruz & \texttt{Selector*}, \texttt{Click} & Navega por el menú y confirma la opción seleccionada. \\
\hline
\end{tabular}
\end{table}

El comando de garra está enlazado actualmente a \textbf{R3}
(\texttt{rightStickPress}); no se utilizan R2/L2 para esta función.

El perfil VR--2 utiliza \texttt{RobotBasic.inputactions}. El dispositivo se
integra como un HID de GrupoVRobotics, con el nombre de producto
\texttt{Joystick VR-v2}. Sus comandos se resumen en la
Tabla~\ref{tab:vr2_commands}.

\begin{table}[H]
\centering
\small
\setlength{\tabcolsep}{4pt}
\renewcommand{\arraystretch}{1.18}
\caption{Comandos implementados para el joystick VR--2 con RPi Pico.}
\label{tab:vr2_commands}
\begin{tabular}{p{3.2cm}p{3.5cm}p{8.9cm}}
\hline
\textbf{Control físico} & \textbf{Acción de entrada} & \textbf{Función} \\
\hline
Palanca principal, eje 1 & \texttt{Mover X} & Traslación lateral del TCP; avance/retroceso de cámara; eje horizontal del menú. \\
Palanca principal, eje 2 & \texttt{Mover Z} & Traslación frontal del TCP; desplazamiento lateral de cámara. Junto con \texttt{Mover X} genera el gesto circular del modo J6. \\
Palanca de altura & \texttt{Mover Y} & Traslación vertical del TCP; eje vertical del menú. \\
Trigger Gripper, clic corto & \texttt{Gripper} & Alterna la apertura/cierre de la garra al soltar el pulsador. \\
Trigger Gripper, mantener \SI{1}{\second} & \texttt{Gripper} & Ejecuta el retorno de J6 a referencia sin alternar la garra. \\
Botón Menú/Toggle, clic corto & \texttt{Camera Toggle} & Conmuta entre Robot y Cámara; cuando el menú está abierto actúa como confirmación. \\
Botón Menú/Toggle, doble clic & \texttt{Camera Toggle} & Activa o desactiva el modo J6 exclusivo. \\
Botón Menú/Toggle, mantener \SI{2}{\second} & \texttt{Camera Toggle} & Abre o cierra el menú de pausa. \\
\hline
\end{tabular}
\end{table}

En el perfil VR--2 la cámara admite traslación con la palanca principal. No se
asignan acciones \texttt{ViewUp}/\texttt{ViewSide}, por lo que su orientación no
se modifica desde este dispositivo en la configuración actual. Los comandos
directos L1/R1 y Cuadrado son propios del perfil PS4; en VR--2, J6 se controla con
la palanca principal y se referencia manteniendo el trigger de la garra.

\subsection{Modos de operación y gestos del botón de cámara}

El componente \texttt{JoystickAdapter} publica el indicador global
\texttt{IsCameraMode} para el modo cámara. También utiliza
\texttt{IsJ6ExclusiveMode} para el modo J6 exclusivo. El control distingue los
siguientes estados:

\begin{itemize}
    \item \textbf{Modo Robot}: los ejes del joystick mueven el
          \emph{Tool Center Point} (TCP) del manipulador.
    \item \textbf{Modo Cámara}: se detiene el brazo y los ejes controlan la
          cámara principal.
    \item \textbf{Modo J6 exclusivo}: se detiene la traslación del TCP y las
          entradas se destinan únicamente al giro de J6.
    \item \textbf{Pausa/menú}: se suprimen los comandos de brazo, cámara y garra
          mientras se navega por las opciones.
\end{itemize}

El mismo botón distingue tres gestos: un clic corto conmuta Robot/Cámara después
de una ventana de \SI{0.4}{\second}; un doble clic activa o desactiva el modo J6;
y una pulsación de \SI{2}{\second} abre o cierra el menú de pausa. En PS4 este
botón es L3 y en VR--2 es el botón Menú/Toggle.

\subsection{Componente \texttt{JoystickAdapter}}

\texttt{JoystickAdapter} se adjunta al objeto raíz del robot (p.ej.\ \emph{KUKA
KR210 R3100-2}) y referencia al \texttt{Controller} provisto por la librería
\emph{Flange}. En cada \texttt{FixedUpdate} (cuando \texttt{IsCameraMode = false})
se construye el incremento de posición en espacio mundo:

\begin{equation}
    \Delta\mathbf{p} = \bigl(\hat{\mathbf{d}}_X \cdot \sigma_X \cdot u_X
                           + \hat{\mathbf{Y}}_{\text{mundo}} \cdot u_Y
                           + \hat{\mathbf{d}}_Z \cdot \sigma_Z \cdot u_Z\bigr)
                       \cdot v \cdot \Delta t
    \label{eq:delta_p}
\end{equation}

donde $u_X, u_Y, u_Z \in [-1,1]$ son los valores leídos de las acciones
\texttt{MoveX}, \texttt{MoveY} y \texttt{MoveZ}; $v$ es la velocidad máxima
configurada en el Inspector; $\hat{\mathbf{d}}_X$ y $\hat{\mathbf{d}}_Z$ son los
versores mundo asignados dinámicamente a cada palanca (ver
Sección~\ref{sec:remap}); y $\sigma_X, \sigma_Z \in \{-1, +1\}$ son los signos
de sentido positivo resultantes del remapeo geométrico.

El incremento se aplica sobre la pose actual del TCP obtenida del observador de
Flange y se pasa al solver de cinemática inversa. Si la solución es válida, se
actualizan los valores articulares del grupo mecánico.

\subsection{Modo J6 exclusivo}

El modo J6 exclusivo se activa o desactiva con doble clic sobre el comando de
cámara. Durante este modo se bloquea el movimiento cartesiano del TCP y se
conservan J1--J5. El operador dispone de los siguientes comandos:

\begin{itemize}
    \item El movimiento circular del stick izquierdo de PS4 o de la palanca
          principal de VR--2 gira J6 en forma proporcional al recorrido realizado.
    \item En PS4, L1 gira en sentido antihorario y R1 en sentido horario.
    \item Al centrar la palanca, J6 se detiene en la posición alcanzada.
    \item Cuadrado en PS4 inicia el retorno de J6 a su referencia física.
    \item Mantener el comando de garra durante \SI{1}{\second} ejecuta el mismo
          retorno tanto en PS4 como en VR--2.
\end{itemize}

Mientras el modo está activo se muestra un dial con el ángulo actual de J6. Al
salir se recupera el control cartesiano normal del TCP.

\subsection{Control de la garra}

Una pulsación corta del comando \texttt{Gripper} alterna apertura y cierre. La
orden se ejecuta al soltar el pulsador, lo que permite distinguirla de la
pulsación prolongada utilizada para referenciar J6. En PS4 se utiliza R3 y en
VR--2 el Trigger Gripper.

\subsection{Componente \texttt{CameraJoystickController}}

\texttt{CameraJoystickController} se adjunta a la \emph{Main Camera} y solo
actúa cuando \texttt{IsCameraMode = true}. Implementa un movimiento estilo
\emph{first-person shooter}:

\begin{itemize}
    \item \textbf{Traslación}: el analógico izquierdo genera desplazamiento
          relativo a la orientación actual de la cámara (avance/retroceso con
          \texttt{MoveForward}, lateral con \texttt{MoveSide}).
    \item \textbf{Rotación}: el analógico derecho controla el \emph{yaw}
          (rotación en $Y$ global, acción \texttt{ViewSide}) y el \emph{pitch}
          (rotación en $X$ local, acción \texttt{ViewUp}), con un límite
          configurable de $\pm 80^{\circ}$ para evitar el volteo de la cámara.
\end{itemize}

En VR--2 se encuentran asignadas las acciones de traslación de cámara, pero no las
de rotación. Al regresar al modo Robot se aplica el remapeo geométrico descrito en
la Sección~\ref{sec:remap}.

\subsection{Pausa, menú, selección de perfil y calibración}

El menú puede abrirse con Options en PS4 o manteniendo \SI{2}{\second} el comando
de cámara en ambos perfiles. Mientras está visible, el sistema suprime las
entradas de movimiento y garra. En PS4 se navega con la cruceta y se confirma con
Cruz; en VR--2 se navega con la palanca principal y la palanca de altura y se
confirma con el botón Menú/Toggle.

Desde el menú se puede seleccionar PS4 o VR--2, calibrar o restablecer los ejes
analógicos de VR--2 y alternar el modo de orientación del TCP. La calibración no
se ofrece para PS4. Al cerrar el menú se exige liberar el botón de cámara antes de
admitir una nueva orden, evitando una conmutación involuntaria.

% ─── Remapeo geométrico ───────────────────────────────────────────────────────
\subsection{Remapeo geométrico de direcciones al retomar el control del robot}
\label{sec:remap}

Cuando el operador navega libremente por la escena y luego retoma el control del
robot (transición \texttt{IsCameraMode}: \texttt{true} $\to$ \texttt{false}), la
dirección ``al frente'' del punto de vista puede no coincidir con los ejes locales
del efector final. Para que los analógicos siempre empujen el TCP en la dirección
intuitiva desde la posición de la cámara, se desarrolló un método geométrico
basado en proyecciones en el plano horizontal $XZ$ del mundo.

\subsubsection*{Paso 1 --- Vector de dirección cámara--efector}

Se calcula el vector $\overrightarrow{C}$ como la diferencia entre la posición
del TCP ($\mathbf{A}$) y la posición de la cámara ($\mathbf{B}$), proyectado
en el plano horizontal (componente $Y = 0$) y normalizado:

\begin{equation}
    \overrightarrow{C} = \frac{(\mathbf{A} - \mathbf{B})\big|_{Y=0}}
                              {\|(\mathbf{A} - \mathbf{B})\big|_{Y=0}\|}
    \label{eq:vecC}
\end{equation}

\subsubsection*{Paso 2 --- Proyección de los ejes locales del efector}

Los versores \emph{forward} ($\hat{\mathbf{f}}$) y \emph{right} ($\hat{\mathbf{r}}$)
del TCP se proyectan también sobre el plano $XZ$ del mundo:

\begin{equation}
    \hat{\mathbf{f}}_{XZ} = \frac{\hat{\mathbf{f}}\big|_{Y=0}}{\|\hat{\mathbf{f}}\big|_{Y=0}\|},
    \qquad
    \hat{\mathbf{r}}_{XZ} = \frac{\hat{\mathbf{r}}\big|_{Y=0}}{\|\hat{\mathbf{r}}\big|_{Y=0}\|}
    \label{eq:proj_axes}
\end{equation}

\subsubsection*{Paso 3 --- Determinación del eje más alineado con $\overrightarrow{C}$}

Se calculan los ángulos entre $\overrightarrow{C}$ y cada eje proyectado mediante
el producto escalar:

\begin{equation}
    \alpha_f = \angle(\overrightarrow{C},\, \hat{\mathbf{f}}_{XZ}), \qquad
    \alpha_r = \angle(\overrightarrow{C},\, \hat{\mathbf{r}}_{XZ})
    \label{eq:angles}
\end{equation}

Dado que \texttt{Vector3.Angle} devuelve valores en $[0^{\circ}, 180^{\circ}]$, se ajustan al
rango $[-90^{\circ}, 90^{\circ}]$ de modo que un ángulo negativo indique que el eje físico
apunta en sentido opuesto a $\overrightarrow{C}$:

\begin{equation}
    \alpha'_i =
    \begin{cases}
        \alpha_i - 180^{\circ} & \text{si } \alpha_i > 90^{\circ} \\
        \alpha_i        & \text{si } \alpha_i \leq 90^{\circ}
    \end{cases}
    \label{eq:adj_angles}
\end{equation}

El eje con menor $|\alpha'_i|$ es el más alineado con $\overrightarrow{C}$.

\subsubsection*{Paso 4 --- Asignación de \texttt{MoveZ} y determinación de signo}

El eje más alineado con $\overrightarrow{C}$ se asigna a la palanca
\texttt{MoveZ} (avance/retroceso). Su signo positivo se determina por el ángulo
original:

\begin{equation}
    \sigma_Z =
    \begin{cases}
        -1 & \text{si } \alpha_Z > 90^{\circ} \quad \text{(el eje apunta al lado opuesto de } \overrightarrow{C}\text{)} \\
        +1 & \text{si } \alpha_Z \leq 90^{\circ}
    \end{cases}
    \label{eq:signZ}
\end{equation}

\subsubsection*{Paso 5 --- Asignación de \texttt{MoveX} y determinación de signo}

El eje restante se asigna a \texttt{MoveX}. Su signo positivo se determina
evaluando si el producto vectorial entre dicho eje ($\hat{\mathbf{e}}_X$) y
$\overrightarrow{C}$ es paralelo al versor $+Y$ del mundo:

\begin{equation}
    \sigma_X =
    \begin{cases}
        +1 & \text{si } \bigl(\hat{\mathbf{e}}_X \times \overrightarrow{C}\bigr) \cdot \hat{\mathbf{Y}} \geq 0 \\
        -1 & \text{en caso contrario}
    \end{cases}
    \label{eq:signX}
\end{equation}

De este modo, empujar el analógico izquierdo hacia la derecha siempre desplaza
el TCP hacia la derecha desde la perspectiva del operador, independientemente de
la posición de la cámara en la escena.

\subsubsection*{Ejemplo numérico}

Tomando como referencia el diagrama de la Figura~\ref{fig:desarrollo},
se consideran los valores:

\begin{align*}
    \mathbf{A} &= (0.2,\; 1.9) \quad \text{(TCP en plano XZ)} \\
    \mathbf{B} &= (2,\; 2.55)  \quad \text{(cámara en plano XZ)} \\
    \overrightarrow{C} &= \mathbf{A} - \mathbf{B} = (-1.8,\; -0.65)
\end{align*}

Los ángulos calculados resultan:

\begin{align*}
    \alpha_f &= 160.1^{\circ}  \quad \Rightarrow \quad \alpha'_f = -19.9^{\circ} \\
    \alpha_r &= 109.86^{\circ} \quad \Rightarrow \quad \alpha'_r = -70.14^{\circ}
\end{align*}

Como $|\alpha'_f| < |\alpha'_r|$, el eje \emph{forward} del TCP se asigna a
\texttt{MoveZ}. Dado que $\alpha_f = 160.1^{\circ} > 90^{\circ}$, se invierte el sentido:
$\sigma_Z = -1$. El eje \emph{right} se asigna a \texttt{MoveX} y su signo se
determina por el producto vectorial según la ecuación~\eqref{eq:signX}.

\begin{figure}[H]
    \centering
    \includegraphics[width=0.65\textwidth]{figura_remapeo}
    \caption{Plano horizontal $XZ$ del mundo mostrando la posición de la cámara
             (MC), el TCP del manipulador (KUKA), el vector $\overrightarrow{C}$
             y los versores $\hat{X}_A$ y $\hat{Z}_A$ del efector final.
             Generado con el script de MATLAB adjunto.}
    \label{fig:desarrollo}
\end{figure}
